%%%%%%%%%%  scan idx 49  |  folio 36  %%%%%%%%%%
ment de $K$ tel que $\lambda^iN=0$ pour $i$ assez grand. On peut alors, pour tout $x\in K$ et tout entier $p\geq1$, considérer les éléments $\lambda^p(N,x)$, qui s'expriment par des polynômes universels à coefficients entiers en les $\lambda^iN$ et les $\lambda^jx$. On a $\lambda^1(N,x)=x$ et l'on a la formule (1.32) \emph{dans} $K$. Pour exprimer les propriétés des opérations $x\mapsto\lambda^p(N,x)$, il est commode d'introduite un nouvel $\lambda$-anneau noté $K_N$. En tant qu'anneau commutatif $K_N$ s'obtient par adjonction d'une unité à l'anneau dont le groupe additif sous-jacent est $K$ et dont la multiplication est donnée par
\begin{equation}\tag{1.34}
x_Ny=xy\mu\qquad (\mu=\lambda_{-1}(N)).
\end{equation}
L'application $\lambda_t$ de $K_N$ dans $K_N[[t]]$ est définie par :
\begin{equation}\tag{1.35}
\lambda_t(n,x)=(1+t)^n\sum_{p\geq0}\lambda^p(N,x)t^p,
\end{equation}
où l'on a posé $\lambda^0(N,x)=(1,0)$ (élément unité de $K_N$). Je dis que, \emph{muni de ces opérations, $K_N$ est un $\lambda$-anneau spécial}. Il s'agit donc de vérifier les formules $\lambda_t(X+Y)=\lambda_t(X)\lambda_t(Y)$, $\lambda_t(XY)=\lambda_t(X).\lambda_t(Y)$, et $\lambda_t(\lambda^i(X))=\lambda^i(\lambda_t(X))$ pour $X,Y\in K_N$. On peut se borner à vérifier ces formules lorsque $X$ et $Y$ sont dans $K$, et l'on constate qu'il suffit de le vérifier lorsque $K$ est le $\lambda$-anneau spécial libre engendré par $X,Y$ et $N$ soumis aux relations $\lambda^iN=0$ pour $i>q$. Mais l'homomorphisme de groupes
\begin{equation}\tag{1.36}
K_N\longrightarrow K,
\end{equation}
qui transforme unité en unité et sur $K$ se réduit à l'homomorphisme $x\longrightarrow x\mu$, est un homomorphisme compatible avec les structures d'anneaux et les applications $\lambda^i$, comme on le vérifie aussitôt. Dans le cas actuel, 
